\documentclass[10pt]{extarticle}
\usepackage[a4paper,margin=1.5cm]{geometry}
\usepackage[T1]{fontenc}
\usepackage{amsmath,amssymb}
\usepackage{graphicx}
\usepackage[export]{adjustbox}
\usepackage{multicol}
\usepackage{enumitem}
\usepackage{booktabs,array,tabularx}
\setlength{\columnsep}{0.7cm}
\setlength{\parindent}{0pt}
\setlength{\parskip}{3pt}
\pagestyle{empty}
\makeatletter
\newcommand{\qfigdef}[4]{\@namedef{qf@file@#1}{#2}\@namedef{qf@w@#1}{#3}\@namedef{qf@wide@#1}{#4}}
\newcommand{\qf@check}[1]{\@ifundefined{qf@file@#1}{\PackageError{score2}{Unknown figure id #1}{}}{}}
\newcommand{\qf@img}[4][]{\includegraphics[width=#2\dimexpr\@nameuse{qf@w@#4}pt\relax,max width=#3,max height=0.45\textheight,#1]{\@nameuse{qf@file@#4}}}
\newif\ifqf@full
\newif\ifqf@boxed
\newsavebox{\qf@box}
\newcommand{\qf@open}{\global\qf@fulltrue\par\global\qf@boxedtrue
  \setbox\qf@box\vbox\bgroup\hsize\textwidth\linewidth\textwidth}
\newcommand{\qf@emit}[1]{\par\egroup\global\qf@boxedfalse
  \ifdim\dimexpr\ht\qf@box+\dp\qf@box\relax<#1\textheight\noindent\usebox{\qf@box}\par\else\unvbox\qf@box\fi}
\newcommand{\qfullsplit}{\ifqf@boxed\qf@emit{0.85}\fi}
\newcommand{\qf@close}{\ifqf@boxed\qf@emit{0.45}\fi\global\qf@fullfalse}
\newcommand{\qfullbegin}{\end{multicols}\qf@open}
\newcommand{\qfullend}{\qf@close\begin{multicols}{2}}
\newcommand{\qfullstart}{\qf@open}
\newcommand{\qfulldone}{\qf@close}
\newcommand{\qfig}[2][1]{\qf@check{#2}%
  \ifqf@full\qf@blk{#1}{#2}\else
  \ifnum\@nameuse{qf@wide@#2}=1 \expandafter\@firstoftwo\else\expandafter\@secondoftwo\fi
  {\end{multicols}\begin{center}\qf@img{#1}{\textwidth}{#2}\end{center}\begin{multicols}{2}}%
  {\qf@blk{#1}{#2}}\fi}
\newcommand{\qf@blk}[2]{\par\smallskip{\centering\qf@img{#1}{\linewidth}{#2}\par}\smallskip}
\newcommand{\qfigi}[2][1]{\qf@check{#2}\qf@img[valign=c]{#1}{\linewidth}{#2}}
\makeatother
\newcommand{\papersec}[1]{\par\bigskip{\centering\Large\bfseries PAPER-#1\par}\medskip}
\newcommand{\partsec}[2]{\par\medskip\noindent{\large\bfseries PART-#1: #2}\par\smallskip}
\newcommand{\sect}[1]{\par\smallskip\noindent\textbf{#1}\par}
\newcommand{\qn}[1]{\par\smallskip\noindent\textbf{#1.}~}
\newcommand{\paraq}[2]{\par\smallskip\noindent\textbf{Paragraph for Questions #1 and #2}\par}
\newenvironment{qoptions}{\begin{itemize}[leftmargin=1.8em,itemsep=1pt,topsep=2pt,label={}]}{\end{itemize}}
\newcommand{\opt}[2]{\item[(#1)] #2}
% figure registry (generated by score2_tools.py finalize)
\qfigdef{p1_phy_s1_q01_1}{SRG (Enthusiast) - Test 30 (22 Oct 2023)-img/p1_phy_s1_q01_1.png}{79.1}{0}
\qfigdef{p1_phy_s1_q02_1}{SRG (Enthusiast) - Test 30 (22 Oct 2023)-img/p1_phy_s1_q02_1.png}{182.5}{0}
\qfigdef{p1_phy_s1_q03_1}{SRG (Enthusiast) - Test 30 (22 Oct 2023)-img/p1_phy_s1_q03_1.png}{214.3}{0}
\qfigdef{p1_phy_s1_q04_1}{SRG (Enthusiast) - Test 30 (22 Oct 2023)-img/p1_phy_s1_q04_1.png}{200.9}{0}
\qfigdef{p1_phy_s1_q05_1}{SRG (Enthusiast) - Test 30 (22 Oct 2023)-img/p1_phy_s1_q05_1.png}{183.2}{0}
\qfigdef{p1_phy_s1_q07_1}{SRG (Enthusiast) - Test 30 (22 Oct 2023)-img/p1_phy_s1_q07_1.png}{189.6}{0}
\qfigdef{p1_phy_s2_q01_1}{SRG (Enthusiast) - Test 30 (22 Oct 2023)-img/p1_phy_s2_q01_1.png}{130.8}{0}
\qfigdef{p1_phy_s2_q02_1}{SRG (Enthusiast) - Test 30 (22 Oct 2023)-img/p1_phy_s2_q02_1.png}{131.8}{0}
\qfigdef{p1_phy_s2_q03_1}{SRG (Enthusiast) - Test 30 (22 Oct 2023)-img/p1_phy_s2_q03_1.png}{154.8}{0}
\qfigdef{p1_phy_s2_q06_1}{SRG (Enthusiast) - Test 30 (22 Oct 2023)-img/p1_phy_s2_q06_1.png}{157.3}{0}
\qfigdef{p1_phy_s3_q01_1}{SRG (Enthusiast) - Test 30 (22 Oct 2023)-img/p1_phy_s3_q01_1.png}{117.6}{0}
\qfigdef{p1_chem_s1_q01_1}{SRG (Enthusiast) - Test 30 (22 Oct 2023)-img/p1_chem_s1_q01_1.png}{125.1}{0}
\qfigdef{p1_chem_s1_q02_1}{SRG (Enthusiast) - Test 30 (22 Oct 2023)-img/p1_chem_s1_q02_1.png}{99.9}{0}
\qfigdef{p1_chem_s1_q03_1}{SRG (Enthusiast) - Test 30 (22 Oct 2023)-img/p1_chem_s1_q03_1.png}{93.8}{0}
\qfigdef{p1_chem_s1_q03_2}{SRG (Enthusiast) - Test 30 (22 Oct 2023)-img/p1_chem_s1_q03_2.png}{92.5}{0}
\qfigdef{p1_chem_s1_q03_3}{SRG (Enthusiast) - Test 30 (22 Oct 2023)-img/p1_chem_s1_q03_3.png}{94.5}{0}
\qfigdef{p1_chem_s1_q03_4}{SRG (Enthusiast) - Test 30 (22 Oct 2023)-img/p1_chem_s1_q03_4.png}{135.0}{0}
\qfigdef{p1_chem_s1_q08_1}{SRG (Enthusiast) - Test 30 (22 Oct 2023)-img/p1_chem_s1_q08_1.png}{154.3}{0}
\qfigdef{p1_chem_s2_q03_1}{SRG (Enthusiast) - Test 30 (22 Oct 2023)-img/p1_chem_s2_q03_1.png}{230.8}{0}
\qfigdef{p1_chem_s2_q08_1}{SRG (Enthusiast) - Test 30 (22 Oct 2023)-img/p1_chem_s2_q08_1.png}{104.8}{0}
\qfigdef{p2_phy_s1_q01_1}{SRG (Enthusiast) - Test 30 (22 Oct 2023)-img/p2_phy_s1_q01_1.png}{101.5}{0}
\qfigdef{p2_phy_s1_q06_1}{SRG (Enthusiast) - Test 30 (22 Oct 2023)-img/p2_phy_s1_q06_1.png}{155.8}{0}
\qfigdef{p2_phy_s1_q07_1}{SRG (Enthusiast) - Test 30 (22 Oct 2023)-img/p2_phy_s1_q07_1.png}{133.2}{0}
\qfigdef{p2_phy_s2_para11_12_1}{SRG (Enthusiast) - Test 30 (22 Oct 2023)-img/p2_phy_s2_para11_12_1.png}{128.6}{0}
\qfigdef{p2_phy_s3_q07_1}{SRG (Enthusiast) - Test 30 (22 Oct 2023)-img/p2_phy_s3_q07_1.png}{134.0}{0}
\qfigdef{p2_phy_s3_q08_1}{SRG (Enthusiast) - Test 30 (22 Oct 2023)-img/p2_phy_s3_q08_1.png}{173.0}{0}
\qfigdef{p2_chem_s1_q06_1}{SRG (Enthusiast) - Test 30 (22 Oct 2023)-img/p2_chem_s1_q06_1.png}{86.9}{0}
\qfigdef{p2_chem_s1_q07_1}{SRG (Enthusiast) - Test 30 (22 Oct 2023)-img/p2_chem_s1_q07_1.png}{61.3}{0}
\qfigdef{p2_chem_s1_q08_1}{SRG (Enthusiast) - Test 30 (22 Oct 2023)-img/p2_chem_s1_q08_1.png}{65.8}{0}

\begin{document}
\begin{center}{\Large\bfseries SRG (Enthusiast) --- Test 30 (22 Oct 2023)}\end{center}
\begin{multicols}{2}
\papersec{1}
\partsec{1}{PHYSICS}
\sect{Section-I}
\qn{1} A loop PQR formed by three identical uniform conducting rods each of length `$a$' is suspended from one of its vertices (P) so that it can rotate about horizontal fixed smooth axis CD. Initially plane of loop is in vertical plane. A constant current `$i$' is flowing in the loop. Total mass of the loop is `$m$'. At $t=0$, a uniform magnetic field of strength $B$ directed vertically upwards is switched on. Acceleration due to gravity is `$g$'.
\qfig{p1_phy_s1_q01_1}
\begin{qoptions}
\opt{A}{Minimum value of $B$ so that the plane of the loop becomes horizontal (even for an instant) during its subsequent motion is $\dfrac{2mg}{3ia}$}
\opt{B}{Minimum value of $B$ so that the plane of the loop becomes horizontal (even for an instant) during its subsequent motion is $\dfrac{4mg}{3ia}$}
\opt{C}{The equilibrium position (angle of the plane of the loop with vertical) is given by $\tan^{-1}\left(\dfrac{3Bai}{2mg}\right)$}
\opt{D}{The equilibrium position (angle of the plane of the loop with vertical) is given by $\tan^{-1}\left(\dfrac{3Bai}{4mg}\right)$}
\end{qoptions}
\qn{2} Two long conducting rectangular strips are placed parallel to each other at a distance $d$ as shown. Strip A is connected to a non-conducting spring of spring constant $k$. Strip B is fixed. Given $\ell\gg b\gg d$. A current $i$ flows uniformly along the length $\ell$ of each strip and the currents in the strips are anti-parallel to each other. The initial deformation in the spring keeps the system in static equilibrium. Now at $t=0$, the currents in both strips are halved and made parallel to each other suddenly. (Neglect gravity and assume force on A is due to B and spring only and consider A and B as large sheets and neglect end effects)
\qfig{p1_phy_s1_q02_1}
\begin{qoptions}
\opt{A}{The instantaneous force on strip A at $t=0^+$ is given by $\dfrac{5\mu_0i^2\ell}{8b}$}
\opt{B}{The instantaneous force on strip A at $t=0^+$ is given by $\dfrac{3\mu_0i^2\ell}{8b}$}
\opt{C}{Elongation in spring at the instant $t=0^+$ is given by $\dfrac{\mu_0i^2\ell}{2bk}$}
\opt{D}{Elongation in spring at the instant $t=0^+$ is given by $\dfrac{\mu_0i^2\ell}{4bk}$}
\end{qoptions}
\qn{3} The figure-2, figure-3, figure-4 and figure-5 represents the graphs of the source voltage $V$, resistor voltage $V_R$, inductor voltage $V_L$ and capacitor voltage $V_C$ as a function of time $t$ of the series circuit with impedance $500\,\Omega$ as shown in the figure-1 $\left(t_0=\dfrac{\pi}{20}\,\text{ms}\right)$. Choose correct option(s).
\qfig{p1_phy_s1_q03_1}
\begin{qoptions}
\opt{A}{$R=400\,\Omega$, $C=0.5\,\mu$F and $L=50$ mH}
\opt{B}{$R=300\,\Omega$, $C=0.5\,\mu$F and $L=60$ mH}
\opt{C}{$V_{10}=30$ V}
\opt{D}{$I=0.1\,\mathrm{A}\cos(10000\,t)$ Amp and current lag the source voltage}
\end{qoptions}
\qn{4} In given young's double slit experiment a point sources of wavelength $\lambda=500$ nm is placed slightly off the central axis. Point A and B are in front of upper and lower slit respectively. If intensity from upper \& lower slits is $I_0$ each on the screen, then
\qfig{p1_phy_s1_q04_1}
\begin{qoptions}
\opt{A}{Intensity at A, O and B is $4I_0$, $2I_0$ and $4I_0$ respectively}
\opt{B}{70th maxima is formed at B}
\opt{C}{30th maxima is formed at A}
\opt{D}{If a glass slab of thickness 0.02 mm and refractive index $\mu=\dfrac32$ is kept in front of upper slit then central maxima formed exactly at O.}
\end{qoptions}
\qn{5} Two sound source A and B start moving on circular path of radius $R$ with angular speed $\omega$ and $5\omega$ respectively. An observer is at position `O' on the line joining the centers. At $t=0$ position of sources are represented in the figure. [Speed of sound is 320 m/s, $R=20$ m, $\omega=\dfrac15$ rad/s]
\qfig{p1_phy_s1_q05_1}
\begin{qoptions}
\opt{A}{Maximum beat frequency observed by observer is 7.48 Hz}
\opt{B}{Maximum beat frequency observed by observer is 9.48 Hz}
\opt{C}{Minimum beat frequency observed by observer is 6.28 Hz}
\opt{D}{Minimum beat frequency observed by observer is 3.14 Hz}
\end{qoptions}
\qn{6} In an AC series circuit, the instantaneous current is zero when the instantaneous voltage is maximum. Connected to the source may be a
\begin{qoptions}
\opt{A}{Pure inductor}
\opt{B}{Pure capacitor}
\opt{C}{Pure resistor}
\opt{D}{Combination of an inductor and capacitor}
\end{qoptions}
\qn{7} Following plots show Magnetization ($M$) vs magnetising field ($H$) and Magnetic susceptibility ($\chi$) vs temperature ($T$) graph :
\qfig{p1_phy_s1_q07_1}
Which of the following combination will be represented by a diamagnetic material?
\begin{qoptions}
\opt{A}{a}
\opt{B}{b}
\opt{C}{c}
\opt{D}{d}
\end{qoptions}
\qn{8} Angular width of central maximum in the Fraunhofer diffraction pattern of a slit is measured. The slit is illuminated by light of wavelength $6000$\,\AA. When the slit is illuminated by light of another wavelength, the angular width decreases by 30\%. Calculate the wavelength of this light. The same decrease in the angular width of central maximum is obtained when the original apparatus is immersed in a liquid. Find refractive index of the liquid.
\begin{qoptions}
\opt{A}{$\lambda=4200$\,\AA}
\opt{B}{$\mu=1.43$}
\opt{C}{$\lambda=8400$\,\AA}
\opt{D}{$\mu=2.46$}
\end{qoptions}
\sect{Section-II}
\qn{1} Two straight infinitely long current wires lying along the $x$-axis and the $y$-axis carry currents 100 A and 200 A as indicated in the adjoining figure. A rigid square metallic wire frame ABCD of side 10 cm and resistance ($R=10\,\Omega$) is connected to a capacitor $C_1=1$ mF as shown in the figure. The wire frame, which lies in the $x$-$y$ plane, moves with a constant velocity $30\left(\hat i+\hat j\right)$ m/s. At the instant, when the point A is located at (20 cm, 20 cm) the charge $q_0$ on the capacitor (as shown in the figure) is 0.2\,$\mu$C. If, at the above mentioned instant, the current flowing in the circuit is given by $10^{-n}$ amp. Fill the value of $n$ in nearest integer.
\qfig{p1_phy_s2_q01_1}
\qn{2} Two identical beams A and B of plane coherent waves of the same intensity and wavelength $\lambda$ fall on cylindrical screen. The angle between the directions of the beam propagations is $\theta$. Consider a point P on the screen at angular position $\phi$ from the beam A as shown in the figure. Find the distance in nanometers between adjacent interference fringes on the screen near the point P. Assume the distance $\beta$ between adjacent fringes is much less than the radius of the cylinder. (Given : $\theta=60^\circ$, $\phi=15^\circ$, $\lambda\sqrt2$ nm)
\qfig{p1_phy_s2_q02_1}
\qn{3} The sides of a square shaped frame, shown in the diagram are made of wires of resistances $r_1$ and $r_2$. The frame is placed in a uniform magnetic field, which is perpendicular to the plane of figure and is increasing at a constant rate. The resistance that should be connected across the diagonal of frame, so that this wire is warmed up at greatest rate is given by $R=\dfrac{nr_1r_2}{r_1+r_2}$. Find $n$.
\qfig{p1_phy_s2_q03_1}
\qn{4} A choke coil is connected to an AC source. Some power is dissipated in coil. Now a capacitor is connected in series with the coil and source is same, then average power dissipated in the circuit does not change. If ratio of inductive reactance of coil and capacitive reactance of capacitor is $\dfrac1x$ then find $x$
\qn{5} At a particular instant a source of sound of frequency 100 Hz is at (2m, 1m) and an observer is at (5m, 5m). The velocity of source and observer are $15\left(2\hat i+\hat j\right)$ m/s and $5\hat i+15\hat j$ m/s respectively at this instant. The velocity of sound in air is 330 m/s. Find the frequency (in Hz) of the sound received by the observer.
\qn{6} In a particular set up, interference is observed on an screen as shown. Wavelength of light emitted by source is $4000$\,\AA. Out of two coherent sources, one is obtained by reflection with a plane mirror of length 2m. If $I_{\max}$ is the maximum intensity in interference pattern, find number of such points in the pattern.
\qfig{p1_phy_s2_q06_1}
\qn{7} In a uniform magnetic field, the magnetic needle has a magnetic moment $9.85\times10^{-2}\,\mathrm{A/m^2}$ and moment of inertia $5\times10^{-6}$. If it performs 10 complete oscillations in 5 seconds then the magnitude of the magnetic field is \underline{\hspace{1.5cm}} mT. [Take $\pi^2$ as 9.85]
\qn{8} The value of aluminium susceptibility is $2.2\times10^{-5}$. The percentage increase in the magnetic field if space within a current carrying toroid is filled with aluminium is $\dfrac{x}{10^4}$. Then the value of $x$ is .
\sect{Section-IV}
\qn{1} In the figure below two strings of different linear mass density $\mu_1$ and $\mu_2$ and length $\ell_1$ and $\ell_2$ are connected to each other. The string passes over the ideal pulley and connected to weight `W'. When excited with a tuning fork the string vibrates such that the joint between the two wires does not move at all. If $N_1$ and $N_2$ are the number of harmonics in the first and second wire respectively, match the enteries in column-I and column-II.
\qfig{p1_phy_s3_q01_1}
\begin{center}\small\renewcommand{\arraystretch}{2.2}
\begin{tabular}{|c|c||c|l|}\hline
\multicolumn{2}{|c||}{\textbf{Column-I}} & \multicolumn{2}{c|}{\textbf{Column-II}}\\\hline
(A) & $\dfrac{\mu_1}{\mu_2}=4:1$ & (P) & $\ell_1=3\ell_2,\ N_1:N_2=1:1$\\\hline
(B) & $\dfrac{\mu_1}{\mu_2}=1:9$ & (Q) & $\ell_1=\ell_2,\ N_1:N_2=1:2$\\\hline
(C) & $\dfrac{\mu_2}{\mu_1}=1:4$ & (R) & $\ell_1:\ell_2=1:6,\ N_1:N_2=1:3$\\\hline
(D) & $\dfrac{\mu_1}{\mu_2}=4:9$ & (S) & $2\ell_1=\ell_2,\ N_1:N_2=1:3$\\\hline
\end{tabular}
\end{center}
\qn{2} Column I represents the standing waves in air columns and string. Column II represents frequency of the note. Match the column-I with column-II. [$v$ = velocity of the sound in the medium]
\begin{center}\small\renewcommand{\arraystretch}{2.0}\setlength\tabcolsep{3pt}
\begin{tabular}{|c|p{0.40\linewidth}|c|c|}\hline
 & \multicolumn{1}{c|}{\textbf{Column - I}} & & \textbf{Column - II}\\\hline
(A) & Second harmonic for the tube open at both ends & (p) & $\dfrac{v}{4\ell}$\\\hline
(B) & Fundamental frequency for the tube closed at one end & (q) & $\dfrac{v}{2\ell}$\\\hline
(C) & First overtone for the tube closed at one end & (r) & $\dfrac{3v}{4\ell}$\\\hline
(D) & Fundamental frequency for the string fixed at both ends & (s) & $\dfrac{v}{\ell}$\\\hline
\end{tabular}
\end{center}
\partsec{2}{CHEMISTRY}
\sect{Section-I}
\qn{1} Which of the following is \textbf{CORRECT} for the given pair of compounds? [Left molecule is (X), right one, (Y)]
\qfig{p1_chem_s1_q01_1}
\begin{qoptions}
\opt{A}{X and Y are enantiomers}
\opt{B}{X and Y are disatereomers}
\opt{C}{X is a D-sugar and Y is an L-sugar}
\opt{D}{If (X) is laevorotatory, (Y) would be dextorotatory.}
\end{qoptions}
\qn{2} Which of the following statement (s) is/are \textbf{CORRECT}?
\qfig{p1_chem_s1_q02_1}
\begin{qoptions}
\opt{A}{The packing efficiency of the two-dimensional ensemble of 5 atoms, shown in the figure is 78.5\%}
\opt{B}{The packing efficiency of the three-dimensional ensemble of repetition of the given structure in ABAB fashion would be 78.5\%}
\opt{C}{The 2-D ensemble will be classified as square close packing}
\opt{D}{The 3-D ensemble (as described in option-B) will be classified as body centred cubic}
\end{qoptions}
\qn{3} Select the correct HORIZONTAL matches from the following table:
\begin{center}\small\renewcommand{\arraystretch}{1.2}\setlength\tabcolsep{3pt}
\begin{tabular}{|c|p{0.40\linewidth}|p{0.30\linewidth}|}\hline
 & \centering Curve / Diagram & \centering\arraybackslash Characteristic / Result / Inference\\\hline
A & \qfigi{p1_chem_s1_q03_1} & W is the most volatile liquid\\\hline
B & \qfigi{p1_chem_s1_q03_2} & Miscible liquids with maximum boiling azeotrope\\\hline
C & \qfigi{p1_chem_s1_q03_3} & Variation of thermodynamics quantities of ideal binary solution\\\hline
D & \qfigi{p1_chem_s1_q03_4} & Boiling point order of the liquids : P $<$ Q $<$ R\\\hline
\end{tabular}
\end{center}
\qn{4} Select the correct HORIZONTAL matches from the following table:
\begin{center}\small\renewcommand{\arraystretch}{1.3}\setlength\tabcolsep{3pt}
\begin{tabular}{|c|p{0.19\linewidth}|p{0.18\linewidth}|p{0.30\linewidth}|}\hline
 & \centering Aq. Solution of & \centering Titrants & \centering\arraybackslash Change of $\lambda$ initially, then after equivalence point [$\lambda$ vs. $V_{\text{titrant}}$ curve]\\\hline
A & Formic acid & Diethyl amine & Increases initially, then becomes constant (doesn't change much)\\\hline
B & $\mathrm{HIO_4}$ & $\mathrm{Ba(OH)_2}$ & Decreases initially, then increases sharply\\\hline
C & $\mathrm{NH_4OH}$ & HBr & Increases initially, then decreases\\\hline
D & $0.1$M $\mathrm{AgNO_3}$ & $1$M KI & Remains constant initially, then decreases\\\hline
\end{tabular}
\end{center}
\qn{5} Two litre solution of a buffer mixture containing $1.0\,\mathrm{M-NaH_2PO_4}$ and $1.0\,\mathrm{M-Na_2HPO_4}$ is placed in two compartments (one litre in each) of an electrolytic cell. The platinum electrodes are inserted in each compartment and 1.25 A current is passed for 965 min. Assuming electrolysis of only water at each compartment, what will be the final pH in each compartment ($\mathrm{p}K_a$ for $\mathrm{H_2PO_4^-}=2.15$, $\log7=0.85$)
\begin{qoptions}
\opt{A}{Anode: 3.00}
\opt{B}{Cathode: 3.00}
\opt{C}{Anode: 1.30}
\opt{D}{Cathode: 1.30}
\end{qoptions}
\qn{6} The polarimeter readings in an experiment of hydrolysis of sucrose are recorded in the following table. Select the \textbf{CORRECT} inference(s) from the given information
\begin{center}\small\setlength\tabcolsep{4pt}
\begin{tabular}{|l|c|c|c|}\hline
Time (min) & 0 & 30 & $\infty$\\\hline
Angle of rotation (degree) & 30 & 20 & $-15$\\\hline
\end{tabular}
\end{center}
Identify the true statement(s)? [$\log2=0.3$, $\log3=0.48$, $\log7=0.84$]
\begin{qoptions}
\opt{A}{The half life of the reaction is 75 min}
\opt{B}{The solution is optically inactive at 120 min}
\opt{C}{The equimolar mixture of the products is dextrorotatory}
\opt{D}{The angle would be $7.5^\circ$ at $t_{\frac12}$.}
\end{qoptions}
\qn{7} Which of the following statement(s) is/are \textbf{TRUE}?
\begin{qoptions}
\opt{A}{Consider the reaction, $\mathrm A\to\mathrm P$ defined to have the rate law: $-\dfrac{d[\mathrm A]}{dt}=k[\mathrm A][\mathrm P]$. The rate of the reaction when plotted vs. time will exhibit a maximum at some time.}
\opt{B}{A positive catalyst in a chemical reaction decreases the $E_{a(\text{forward})}$ and increases the $E_{a(\text{backward})}$.}
\opt{C}{$t_{3/4}$ i.e., time taken to reduce the concentration from $\mathbf a$ to $\dfrac{3a}{4}$ is independent of initial concentration of catalytic decomposition of $\mathrm{NH_3}$ on Pt surface.}
\opt{D}{An organic compound undergoes first-order decomposition and the time taken for its decomposition to $\dfrac18$ and to $\dfrac1{10}$ from its initial concentration are $t_{1/8}$ and $t_{1/10}$. Then value of $\dfrac{10\times t_{1/8}}{t_{1/10}}$ is 9. [$\log_{10}2=0.3$]}
\end{qoptions}
\qn{8} Consider the following methods for preparation of colloids using different methods.
\qfig{p1_chem_s1_q08_1}
Which of the following is \textbf{INCORRECT} for the colloids thus produced?
\begin{qoptions}
\opt{A}{The colloid (A) will be more readily precipitated by $\mathrm{MgCl_2}$ than by $\mathrm{AlCl_3}$}
\opt{B}{The colloid (B) will show cataphoresis followed by coagulation at the destined electrode}
\opt{C}{The colloid (A) with higher zeta potential will be more stable than otherwise}
\opt{D}{The colloid (A) would be effective to treat at patient who had been poisoned by colloidal $\mathrm{As_2O_3}$}
\end{qoptions}
\sect{Section-II}
\qn{1} The total number of oxygen atoms per molecule in each cytosine, thymine and uracil be $\mathbf x$\\
The total number of free hydroxyl groups in glucopyranosyl-fructofuranose be $\mathbf y$\\
The total number of free hydroxyl groups in Methyl 2-deoxyribopyranoside be $\mathbf z$\\
What is $(\mathbf{x+y-z})$
\qn{2} Of the following polymers, $\mathbf{(x)}$ are thermosetting plastics, $\mathbf{(y)}$ are thermoplastics and $\mathbf{(z)}$ are biodegradable polymers. Find the value $(100x+10y+z)$
\begin{center}\small\setlength\tabcolsep{3pt}
\begin{tabular}{|l|l|}\hline
I Polystyrene & II Teflon\\\hline
III Melamine resin & IV Bakelite\\\hline
V PVC & VI Buna-S\\\hline
VII Nylon-66 & VIII Nylon-26\\\hline
IX PHBV & X Gutta percha\\\hline
XI HDPE & XII Polylactic acid\\\hline
\end{tabular}
\end{center}
\qn{3} The number of functional groups in the following compound:
\qfig{p1_chem_s2_q03_1}
Which will react with:\\
(i) Na to evolve $\mathrm{H_2}$ be $\mathbf{(x)}$\\
(ii) $\mathrm{NaBH_4}$ at room temperature and undergo reduction be $\mathbf{(y)}$\\
(iii) And decolourise bromine water be $\mathbf{(z)}$\\
Find $\mathbf{(x+y+z)}$
\qn{4} Iron crystallizes in the two body centered Cubic Lattice the $\alpha-$form below $750^\circ$C and $\gamma-$form above $1400^\circ$C and in a face centered cubic $\beta-$form between those two temperatures. If the maximum ratio of foreign paraticle occupying a hole in the face of a body centered cubic unit cell to the host iron radius is $\dfrac{4\sqrt a-a}{a}$, then value of `$\mathbf a$' is \underline{\hspace{1cm}}
\qn{5} The freezing point of an aqueous solution of KCN containing 0.2 mole/kg water was $-0.80^\circ$C. On adding 0.1 mole of $\mathrm{Cn(CN)_2}$ in the solution containing 1 kg of water, the freezing point of the solution was $-0.6^\circ$C. Assuming that the complex is formed according to the following equation\\
$\mathrm{Cn(CN)_2+mCN\to\left[Cn(CN)_{m+2}\right]^{-m}}$\\
and $\mathrm{Cn(CN)_2}$ is the limiting reactant, the value of $\mathbf m$ is
\qn{6} Stannic chloride is being electolyzed at a power where 100J energy is being consumed every second. 25 batteries with a cell potential of 5V each are connected in series. The volume (in mL) of chlorine evolved per minute at NTP is [molar electronic charge $-96.5$ kilo coulomb] (Closet integer value)
\qn{7} The gaseous decomposition reaction, $\mathrm A(g)\to2\mathrm B(g)+\mathrm C(g)$ is observed to be first order. It is found that after 10 minutes the total pressure of system is 188 torr and after very long time it is 3.88 torr. The rate constant of the reaction (in the unit of minute is)
\qn{8} Accumulation of hydrogen gas over platinum surface obeys some adsorption isotherm, which for a very short interval can be represented as in the diagram below ($x$ and $m$ are plotted for grams while $P$ for atm) The volume (at STP in L) of hydrogen gas that could be absorbed on the surface of 1.67g platinum at a pressure of $8.1\times10^{-3}$ atm is
\qfig{p1_chem_s2_q08_1}
\sect{Section-IV}
\qn{1} Match the following
\begin{center}\small\renewcommand{\arraystretch}{1.3}
\begin{tabular}{|c|l|c|l|}\hline
 & \textbf{Column-I} & & \textbf{Column-II}\\\hline
A & Nylon-6 & P & Condensation\\\hline
B & Silicone & Q & Homopolymer\\\hline
C & Rayon & R & Elastomer\\\hline
D & Buna-N & S & Fiber\\\hline
 & & T & Branched\\\hline
\end{tabular}
\end{center}
\qn{2} Match the following
\begin{center}\small\renewcommand{\arraystretch}{1.3}
\begin{tabular}{|c|l|c|l|}\hline
 & Column-I & & Column-II\\\hline
A & Phenyl Salicylate & P & Antiseptic\\\hline
B & Penicillin & Q & Antibiotic\\\hline
C & Chloramphenicol & R & Bacteriostatic\\\hline
D & Ibuprofen & S & Bactericidal\\\hline
 & & T & Analgesic\\\hline
\end{tabular}
\end{center}
\partsec{3}{MATHEMATICS}
\sect{Section-I}
\qn{1} If $ax^2-bx+c=0$ has two distinct roots lying in the interval $(0,1)$, $a,b,c\in N$, then which of the following may be true
\begin{qoptions}
\opt{A}{$\log_5abc=1$}
\opt{B}{$\log_5abc=2$}
\opt{C}{$\log_5abc=3$}
\opt{D}{$\log_5abc=4$}
\end{qoptions}
\qn{2} The equation $a_8x^8+a_7x^7+a_6x^6+\ldots+a_0=0$ has all its roots positive and real $\left(\text{where }a_8=1,\ a_7=-4,\ a_0=\dfrac1{2^8}\right)$, then
\begin{qoptions}
\opt{A}{$a_1=\dfrac1{2^8}$}
\opt{B}{$a_1=\dfrac{-1}{2^4}$}
\opt{C}{$a_2=\dfrac7{2^4}$}
\opt{D}{$a_3=\dfrac1{2^2}$}
\end{qoptions}
\qn{3} Let $g'(x)>0$ and $f'(x)<0\ \ \forall x\in R$, then
\begin{qoptions}
\opt{A}{$g(f(x+1))>g(f(x-1))$}
\opt{B}{$g(f(x+1))<g(f(x-1))$}
\opt{C}{$f(g(x-1))>f(g(x+1))$}
\opt{D}{$f(g(x-1))<f(g(x+1))$}
\end{qoptions}
\qn{4} For the function $f(x)=x\cos\dfrac1x,\ x\geqslant1$
\begin{qoptions}
\opt{A}{$\displaystyle\lim_{x\to\infty}f'(x)=1$}
\opt{B}{$f'(x)>1$ for all $x\geqslant1$}
\opt{C}{$f''(x)<0$ for all $x\in[1,\infty)$}
\opt{D}{$f''(x)>0$ for all $x\in[1,\infty)$}
\end{qoptions}
\qn{5} The solution of the differential equation $\left(\dfrac{dy}{dx}\right)^2+2y\cot x\dfrac{dy}{dx}=y^2$ is
\begin{qoptions}
\opt{A}{$y-\dfrac{C}{1+\cos x}=0$}
\opt{B}{$y=\dfrac{C}{1-\cos x}$}
\opt{C}{$x=2\sin^{-1}\sqrt{\dfrac{C}{2y}}$}
\opt{D}{$y=2\cos^{-1}\sqrt{\dfrac{C}{2y}}$}
\end{qoptions}
\qn{6} The differential equation of the curve such that the initial ordinate of any tangent is equal to the corresponding subnormal, is
\begin{qoptions}
\opt{A}{A linear differential equation}
\opt{B}{A homogeneous differential equation}
\opt{C}{An differential equation with separable variables}
\opt{D}{None of these}
\end{qoptions}
\qn{7} If $\displaystyle\int\tan^5x\,dx=A\tan^4x+B\tan^2x+g(x)+c$, then
\begin{qoptions}
\opt{A}{$A=\dfrac14,\ B=-\dfrac12$}
\opt{B}{$g(x)=\ln|\sec x|$}
\opt{C}{$g(x)=\ln|\cos x|$}
\opt{D}{$A=-\dfrac14,\ B=\dfrac13$}
\end{qoptions}
\qn{8} If $I=\displaystyle\int_0^1\dfrac{dx}{1+x^{\pi/2}}$, then
\begin{qoptions}
\opt{A}{$I>\ln2$}
\opt{B}{$I<\ln2$}
\opt{C}{$I<\dfrac\pi4$}
\opt{D}{$I>\dfrac\pi4$}
\end{qoptions}
\sect{Section-II}
\qn{1} The value of $\displaystyle\int_{-6}^{6}\max\left(\left|2-|x|\right|,\,4-|x|,\,3\right)dx$ is
\qn{2} If $K\in N$ and $I_k=\displaystyle\int_{-2k\pi}^{2k\pi}|\sin x|\,[\sin x]\,dx$, (where $[.]$ denotes the greatest integer function) then $\displaystyle\sum_{k=1}^{100}I_k$ equals to $m$ then value of $|m|$ is
\qn{3} The quadratic $x^2+ax+b+1=0$ has real roots which are distinct positive integers, then minimum value of $\left(a^2+b^2\right)$ is
\qn{4} If $\alpha,\beta;\ \beta,\gamma$ and $\gamma,\alpha$ are the roots of $a_ix^2+b_ix+c_i=0$ respectively for $i=1,2,3$ (given that $\alpha,\beta,\gamma>0$) respectively if
\[\sum\alpha+\sum\alpha\beta+\alpha\beta\gamma=\left(\prod_{i=1}^{3}\left(\frac{a_i-b_i+c_i}{a_i}\right)\right)^{\frac12}+k\]
, then $|k|$ is equal to \underline{\hspace{1cm}}.
\qn{5} If $f(x)=\begin{vmatrix}x^n&n!&2\\\cos x&\cos\frac{n\pi}2&4\\\sin x&\sin\frac{n\pi}2&8\end{vmatrix}$, then the value of $\dfrac{d^n}{dx^n}[f(x)]$ at $x=0$ is
\qn{6} If $f(x)=\cos\left(\dfrac\pi2[x]-x^3\right)$, $1<x<2$ and $[x]$ = the greatest integer $\leqslant x$, then $f'\left(\sqrt[3]{\dfrac\pi2}\right)=$
\qn{7} If $y^{\frac13}+y^{-\frac13}=2x$ and $\left(x^2-1\right)y''+xy'=ky$, then the value of $k$ is
\qn{8} $f'(x)=\sin x+\sin4x\cdot\cos x$, then $\left(f\left(2x^2+\dfrac\pi2\right)\right)'$ at $x=\sqrt{\dfrac\pi2}$ is $-k\sqrt{2\pi}$, then $k$ is \underline{\hspace{1cm}}.
\sect{Section-IV}
\qn{1} Match the column I and II
\begin{center}\small\renewcommand{\arraystretch}{1.5}\setlength\tabcolsep{3pt}
\begin{tabular}{|c|p{0.40\linewidth}|c|c|}\hline
 & \centering\textbf{Column - I} & & \textbf{Column - II}\\\hline
(A) & Let A be the differential equation of family of lines situated at a constant distance $p$ from the origin, then the degree of differential equation A is & (P) & 0\\\hline
(B) & Value of $\displaystyle\int_{-4}^{-5}e^{(x+5)^2}dx+3\int_{1/3}^{2/3}e^{9\left(x-\frac23\right)^2}dx$ is & (Q) & 1\\\hline
(C) & If each $a_i>0$, then the shortest distance between the point $(0,-3)$ and the curve $y=1+a_1x^2+a_2x^4+\ldots+a_nx^{2n}$ is & (R) & 2\\\hline
(D) & The order of the differential equation of the family of circles with one diameter along the line $y=x$ is & (S) & 4\\\hline
 & & (T) & 3\\\hline
\end{tabular}
\end{center}
\qn{2} Match the column I and II
\begin{center}\footnotesize\renewcommand{\arraystretch}{1.6}\setlength\tabcolsep{3pt}
\begin{tabular}{|c|p{0.40\linewidth}|c|c|}\hline
 & \centering\textbf{Column - I} & & \textbf{Column - II}\\\hline
(A) & The minimum value of $|\sin x+\cos x+$ $\sec x+\operatorname{cosec}x|$ is \underline{\hspace{0.8cm}} & (P) & 0\\\hline
(B) & If the angle of intersection of $x^yy^x=16$ and $x=y$ is $\dfrac\pi k$, then $k$ is \underline{\hspace{0.8cm}} & (Q) & 1\\\hline
(C) & Number of critical points of $f(x)=\dfrac{\left|x^2-1\right|}{x^2}$ is \underline{\hspace{0.8cm}} & (R) & 2\\\hline
(D) & Number of critical points of $\max\{|\sin x|,|\cos x|\}$ in $\left(\dfrac\pi2,\dfrac{3\pi}2\right)$ is/are \underline{\hspace{0.8cm}} & (S) & 3\\\hline
 & & (T) & 4\\\hline
\end{tabular}
\end{center}
\papersec{2}
\partsec{1}{PHYSICS}
\sect{Section-I(i)}
\qn{1} In a series LCR circuit
\qfig{p2_phy_s1_q01_1}
\begin{qoptions}
\opt{A}{The voltage $V_L$ across the inductance leads the current in the circuit by a phase angle of $\dfrac\pi2$}
\opt{B}{The voltage $V_C$ across the capacitance lags behind the current by a phase angle of $\dfrac\pi2$}
\opt{C}{The voltage $V_R$ across the resistance is in phase with the current}
\opt{D}{The voltage across the series combination of L, C and R is $V=V_L-V_C+V_R$}
\end{qoptions}
\qn{2} In a series LCR circuit with an ac source $E_0=50$V, $I_{\rm rms}=0.1$ amp, $R=300\ \Omega$, frequency $\nu=\dfrac{50}{\pi}$Hz. The average electric field energy stored in the capacitor and average magnetic energy stored in the coil are 25 mJ and 5 mJ. Then
\begin{qoptions}
\opt{A}{Capacitance C of capacitor is $20\,\mu$F}
\opt{B}{Inductance L of inductor is 2H}
\opt{C}{Peak voltage of source is 50 V}
\opt{D}{The sum of rms voltage across the three elements is 35.4V}
\end{qoptions}
\qn{3} A sound wave of frequency $f$ travels horizontally to the right. It is reflected from a large vertical plane surface moving to left with a speed $v$. The speed of sound in medium is $C$ :
\begin{qoptions}
\opt{A}{The number of wave striking the surface per second is $f\dfrac{(c+v)}{c}$}
\opt{B}{The wavelength of reflected wave is $\dfrac{c(c-v)}{f(c+v)}$}
\opt{C}{The frequency of the reflected wave is $f\dfrac{(c+v)}{(c-v)}$}
\opt{D}{The number of beats heard by a stationary listener to the left of the reflecting surface is $\dfrac{2vf}{c-v}$}
\end{qoptions}
\qn{4} A transparents slab of thickness `$t$' and refractive index $\mu$ is inserted in front of upper slit of YDSE apparatus. The wavelength of light used is $\lambda$. Assume that there is no absorption of light by the slab. Mark the correct statements.
\begin{qoptions}
\opt{A}{The intensity of dark fringes will be zero, if slits are identical.}
\opt{B}{The change in optical path due to insertion of plate is $\mu t$}
\opt{C}{The change in optical path due to insertion of plate is $(\mu-1)t$}
\opt{D}{For making intensity zero at centre of screen the thickness can be $\dfrac{5\lambda}{2(\mu-1)}$}
\end{qoptions}
\qn{5} In Young's double slit experiment, the two slits are covered by slabs of same thickness but refractive index 1.4 and 1.7. The distance between slits and screen is 1 m and distance between slits is 1 mm and wavelength of coherent source used is $4000$\,\AA\ and the central fringe shifts to the 3rd bright fringe position, then
\begin{qoptions}
\opt{A}{Shift will be towards slab of R.I.-1.7 by 1.2 mm.}
\opt{B}{Shift will be towards slab of R.I.-1.4 by 1.2 mm}
\opt{C}{Slabs are of thickness $4\,\mu$m}
\opt{D}{Slabs are of thickness $4$\,\AA}
\end{qoptions}
\qn{6} A large thin non conducting plate with uniform surface charge density $\sigma$ is moving with constant speed $v$ as shown in the figure. The magnetic field at a small distance above or below the plate in the middle region is
\qfig{p2_phy_s1_q06_1}
\begin{qoptions}
\opt{A}{$\mu_0\sigma v$ in magnitude}
\opt{B}{$\dfrac{\mu_0\sigma v}{2}$ in magnitude}
\opt{C}{Perpendicular to the plate}
\opt{D}{Parallel to plate}
\end{qoptions}
\qn{7} In the circuit shown in the figure (the symbols have their usual meaning), only one of the three switches is kept closed and other two are open during the entire experiment. The value of resistance is $R=20\ \Omega$ When the angular frequency ($\omega$) of the 100 V source is adjusted to 500 rad/s, 1000 rad/s and 2000 rad/s it was found that the current $I$ was 4A, 5A and 4A respectively.
\qfig{p2_phy_s1_q07_1}
\begin{qoptions}
\opt{A}{The switch that was closed is $S_2$}
\opt{B}{The switch that was closed is $S_3$}
\opt{C}{The value of L is 10 mH}
\opt{D}{The value of C is $100\ \mu$F}
\end{qoptions}
\qn{8} Two infinitely long straight wires lie in the $xy$-plane along the lines $x=\pm R$. The wire located at $x=+R$ carries a constant current $I_1$ and the wire located at $x=-R$ carries a constant current $I_2$. A circular loop of radius $R$ is suspended with its centre at $\left(0,0,\sqrt3R\right)$ and in a plane parallel to the $xy$-plane. This loop carries a constant current $I$ in the clockwise direction as seen from above the loop. The current in the wire is taken to be positive if it is in the $+\hat j$ direction. Which of the following statements regarding the magnetic field $\vec B$ is (are) true?
\begin{qoptions}
\opt{A}{If $I_1=I_2$ then $\vec B$ cannot be equal to zero at the origin $(0,0,0)$}
\opt{B}{If $I_1>0$ and $I_2<0$, then $\vec B$ can be equal to zero at the origin $(0,0,0)$}
\opt{C}{If $I_1<0$ and $I_2>0$, then $\vec B$ can be equal to zero at the origin $(0,0,0)$}
\opt{D}{If $I_1=I_2$ then the $z$-component of the magnetic field at the centre of the loop is $\left(-\dfrac{\mu_0I}{2R}\right)$}
\end{qoptions}
\sect{Section-I(ii)}
\paraq{9}{10}
Superposition of waves results in maximum and minimum of intensities such as in case of standing waves. This phenomenon is called as interference. Another type of superposition result in interference in time which is called as beats. In this case waves are analyzed at a fixed point as a function of time. If the two waves are of nearby same frequency are superimposed, at a particular point, intensity of combined waves gives a periodic peak and fall. This phenomenon is beats. If $\omega_1$ and $\omega_2$ are the frequencies of two waves then by superimposed $y=y_1+y_2$ we get at $x=0$,
\[y=\left[2A\cos\left(\frac{\omega_1-\omega_2}{2}\right).t\right]\sin\left[\frac{\omega_1+\omega_2}{2}\right].t\]
Thus amplitude frequency is small and fluctuates slowly. A beat i.e., a maximum of intensity occurs, also intensity depends on square of amplitude. The beat frequency is given by\\
$\omega_{\rm beat}=|\omega_1-\omega_2|$\\
Number of beats per second is called as beat frequency. A normal ear can detect only upto 10 Hz of frequency because of persistence of ear.
\qn{9} If two sound sources of frequency difference 25 Hz are sounded together. Then which of the following is correct ?
\begin{qoptions}
\opt{A}{A normal human ear will hear 25 Hz beat frequency}
\opt{B}{A normal human ear will hear only 10 Hz beat frequency}
\opt{C}{A normal human ear cannot detect this frequency difference}
\opt{D}{A normal human ear can hear maximum of the two frequency sounded together}
\end{qoptions}
\qn{10} The frequency of beats produced in air when two sources of sound are activated, one emitting wavelength 32 cm, other 32.2 cm is (Take $V_{\rm sound}=350$ m/s)
\begin{qoptions}
\opt{A}{7}
\opt{B}{10}
\opt{C}{14}
\opt{D}{18}
\end{qoptions}
\paraq{11}{12}
Two capacitors of capacitance $C$ and $3C$ are charged to potential difference $V_0$ and $2V_0$ respectively, and connected to an inductor of inductance $L$ as shown in figure. Initially the current in the inductor is zero. Now, the switch S is closed.
\qfig{p2_phy_s2_para11_12_1}
\qn{11} The maximum current in the inductor is
\begin{qoptions}
\opt{A}{$\dfrac{3V_0}{2}\sqrt{\dfrac{3C}{L}}$}
\opt{B}{$V_0\sqrt{\dfrac{3C}{L}}$}
\opt{C}{$2V_0\sqrt{\dfrac{3C}{L}}$}
\opt{D}{$V_0\sqrt{\dfrac{C}{L}}$}
\end{qoptions}
\qn{12} Potential difference across capacitor of capacitance $3C$ when the current in the circuit is maximum is:
\begin{qoptions}
\opt{A}{$\dfrac{V_0}{4}$}
\opt{B}{$\dfrac47V_0$}
\opt{C}{$\dfrac54V_0$}
\opt{D}{$\dfrac74V_0$}
\end{qoptions}
\sect{Section-III}
\qn{1} In YDSE, the slits have different widths. As a result, amplitude of waves from two slits are $A$ and $2A$ respectively. If $I_0$ be the maximum intensity of the interference pattern then the intensity of the pattern at a point where the phase difference between waves is $\phi$ is given by $\dfrac{I_0}{P}(5+4\cos\phi)$. Where $P$ is in integer. Find the value of $P$?
\qn{2} A solenoid with inductance $L=7$ mH and resistance $R=44\ \Omega$ is first connected to a source of direct voltage $V_o$ and then to a source of sinusoidal voltage with effective rms value $V=V_o$. At what frequency (in KHz) of the oscillator will the power consumed by the solenoid be $\eta=5.0$ times less than in the former case?
\qn{3} A flood light is covered with a filter that transmits red light. The electric field of the emerging beam is represented by a sinusoidal plane wave\\
$E_x=36\sin\left(1.20\times10^7z-3.6\times10^{15}t\right)$ V/m\\
The average intensity of the light will be $x\dfrac{\mathrm w}{\mathrm m^2}$. Closest integer of $x$ is ?
\qn{4} In YDSE, the intensity of light at a point on the screen is $I$ for a path difference $\lambda$. The intensity of light at a point where the path difference becomes $\dfrac\lambda3$ is $\dfrac IP$. Find the value of $P$?
\qn{5} A magnetic needle pivoted through its centre of mass and is free to rotate in a plane containing uniform magnetic field $200\times10^{-4}$T. When it is displaced slightly from the equilibrium it makes 2 oscillations per second. If the moment of inertia of the needle about the axis of oscillation is $0.75\times10^{-5}$ kg m$^2$ the magnetic moment of the needle is $x$ J/T find $100\,x$ ?
\qn{6} The magnetic moment of magnet of mass 75 gm is $9\times10^{-7}\ \mathrm{Am^2}$. If the density of the material of magnet is $7.5\times10^3\,\dfrac{\mathrm{kg}}{\mathrm{m^3}}$ then intensity of magnetisation will be $x\dfrac{\mathrm A}{\mathrm m}$. Find $100\,x$
\qn{7} Consider a plane inclined at an angle $45^\circ$ with the horizontal has two narrow slits ($S_1$ and $S_2$) separated by a distance $d=\sqrt2$ mm. The screen is placed at a distance of $D=10$ m. A parallel monochromatic light beam of wavelength $5000$\,\AA\ is incident on the slits as shown. If the fringe width of interference pattern on the screen in the region close to the opposite point to the midpoint between the slits is $k\times10^{-3}$ meter, then find the value of $k$.
\qfig{p2_phy_s3_q07_1}
\qn{8} A conducting rod of length `$\ell$' $=2\sqrt5$ meter and mass `$m$' $=4$ kg lies on the horizontal table. Coefficient of friction between the rod and the table is `$\mu$' $=\left(\dfrac12\right)$. If the current in the conductor is 2 A, then find the minimum magnitude of magnetic field strength (in tesla) such that conducting rod just starts to translate along $x$-axis.(taking $g=10$ m/s$^2$) [neglect the radius of rod]
\qfig{p2_phy_s3_q08_1}
\partsec{2}{CHEMISTRY}
\sect{Section-I(i)}
\qn{1} Which of the following is positively charged sols
\begin{qoptions}
\opt{A}{$\mathrm{Al_2O_3.xH_2O}$}
\opt{B}{Gold sols}
\opt{C}{$\mathrm{Sb_2S_3}$}
\opt{D}{Haemoglobin}
\end{qoptions}
\qn{2} For a first order reaction, correct statements are
\begin{qoptions}
\opt{A}{The degree of dissociation is equals to $\left(1-e^{kt}\right)$}
\opt{B}{Half-life period is independent of initial concentrations.}
\opt{C}{The time taken for the completion of 87.5\% reaction is thrice the $t_{1/2}$ of the reaction.}
\opt{D}{Units of rate constant is $\mathrm{sec^{-1}}$.}
\end{qoptions}
\qn{3} The correct statement regarding NaCl crystal ?
\begin{qoptions}
\opt{A}{Second nearest neighbours of $\mathrm{Na^+}$ is 12}
\opt{B}{After removal of all atoms from one body diagonal new formula of NaCl is $\mathrm{Na_4Cl_5}$}
\opt{C}{If radius of $\mathrm{Cl^-}$ is $r_\ominus$ then $4r_\ominus=\sqrt2\,a$}
\opt{D}{If radius of $\mathrm{Na^+}$ is $r_+$ and $\mathrm{Cl^-}$ is $r_-$ then $2r_++2r_-=\dfrac a2$}
\end{qoptions}
\qn{4} Which of the following pair is isotonic in nature ?
\begin{qoptions}
\opt{A}{0.1 M 100 ml $\mathrm{K_4\left[Fe(CN)_6\right]}$ and 0.2 M 50 ml $\mathrm{Al_2(SO_4)_3}$}
\opt{B}{0.1 M 20 ml $\mathrm{Na_2SO_4}$ ($\alpha=75\%$) and 0.2 M 12.5 ml NaCl ($\alpha=100\%$)}
\opt{C}{0.1 M 250 ml $\mathrm{CaCl_2}$ and 0.5 M 150 ml $\mathrm{C_6H_{12}O_6}$}
\opt{D}{0.5 M 500 ml Urea and 0.5 M 250 ml $\mathrm{FeSO_4}$}
\end{qoptions}
\qn{5} Some standard electrode potentials at 298 K are given\\
$\mathrm{Ni^{+2}/Ni}=-0.24$\\
$\mathrm{Cd^{+2}/Cd}=-0.40$\\
$\mathrm{Fe^{+2}/Fe}=-0.44$\\
$\mathrm{Zn^{+2}/Zn}=-0.76$\\
To a solution containing 0.001 M of $\mathrm{X^{+2}}$ and 0.1 M of $\mathrm{Y^{+2}}$, the metal rods X and Y are inserted (at 298K) and connected by a conducting wire. This resulted in dissolution of X. The correct combination(s) of X and Y respectively is/are.
\begin{qoptions}
\opt{A}{Cd and Ni}
\opt{B}{Cd and Fe}
\opt{C}{Ni and Zn}
\opt{D}{Ni and Fe}
\end{qoptions}
\qn{6} The following structure is a
\qfig{p2_chem_s1_q06_1}
\begin{qoptions}
\opt{A}{Aldohexose}
\opt{B}{Aldopentose}
\opt{C}{Pyranose}
\opt{D}{Reducing sugar}
\end{qoptions}
\qn{7} Correct statement for the following polymer
\qfig{p2_chem_s1_q07_1}
\begin{qoptions}
\opt{A}{It is a copolymer between ethene and propene}
\opt{B}{It is a hydrogenated polymer of isoprene}
\opt{C}{It is a Addition polymer}
\opt{D}{It is a condensation polymer}
\end{qoptions}
\qn{8} Correct statement for
\qfig{p2_chem_s1_q08_1}
\begin{qoptions}
\opt{A}{Positive test for $\mathrm{I_2/NaOH}$}
\opt{B}{Positive test for $\mathrm{\left[Ag(NH_3)_2\right]OH}$}
\opt{C}{Positive test for Neutral $\mathrm{FeCl_3}$}
\opt{D}{Positive test for Baeyer's Reagent (Alkaline $\mathrm{KMnO_4}$)}
\end{qoptions}
\sect{Section-I(ii)}
\paraq{9}{10}
The electrochemical cell shown below is a concentration cell $\mathrm{Pb/Pb^{+2}}$ (saturated solution of a sparingly soluble salt, $\mathrm{PbCl_2}$) $\|$ $\mathrm{Pb^{+2}}$ (0.0001 M)/Pb. The emf of the cell depends on the difference in concentrations of $\mathrm{Pb^{+2}}$ ions at the two electrodes. The emf of the cell at 298k is 0.1182V. $\left[\dfrac{2.303R298}{F}=0.0591\right]$
\qn{9} The value of $\Delta G$ (in J mol$^{-1}$) for the given cell is
\begin{qoptions}
\opt{A}{-22812}
\opt{B}{-11406}
\opt{C}{-57000}
\opt{D}{+11406}
\end{qoptions}
\qn{10} The solubility product of $\mathrm{PbCl_2}$ at 298 k based on the information available for the given concentration cell is.
\begin{qoptions}
\opt{A}{$4\times10^{-25}$}
\opt{B}{$10^{-24}$}
\opt{C}{$4\times10^{-24}$}
\opt{D}{$4\times10^{-16}$}
\end{qoptions}
\paraq{11}{12}
In diamond structure, carbon atoms form fcc lattice and 50\% tetrahedral voids occupied by carbon atoms. Every carbon atoms is surrounded tetrahedrally by four carbon atoms with bond length 154 pm. Germanium, silicon and grey tin also crystallise in same way as diamond. ($N_A=6\times10^{23}$)
\qn{11} The mass of diamond unit cell is :
\begin{qoptions}
\opt{A}{96 amu}
\opt{B}{96 gm}
\opt{C}{144 amu}
\opt{D}{144 gm}
\end{qoptions}
\qn{12} The side length of diamond unit cell is (in pm) :
\begin{qoptions}
\opt{A}{154}
\opt{B}{1422.63}
\opt{C}{711.32}
\opt{D}{355.66}
\end{qoptions}
\sect{Section-III}
\qn{1} An element has a F.C.C structure with a cell edge of 500pm. The density of the element is 7.2 g/cm$^3$. Number of atoms are present in 180 gm of the element is $x\times10^{23}$ atoms. The value of $x$ is
\qn{2} A solution of 61 gm of solute (M.wt = 122 gm) in 1000 gm of solvent (molar mass = 92 gm/mole) shows a boiling point elevation of $1.5^\circ$. Assuming that solute is dimerized to the extent of 80\%. Molar enthalpy of vapourisation of solvent is 9.2 kCal/mole. Normal boiling point of solvent is $y\times10^2$k. The value of $y$ is.
\qn{3} Number of reducing sugar from the following D-Glucose, D-Fructose, Sucrose, Maltose, Starch, Cellulose, Amylose, D-Ribose, D-Glyceraldehyde and Lactose
\qn{4} Pentapeptide on hydrolysis gives Asp, Ala, Gly, Lys and Leu. In pentapeptide C-terminal Amino acid should be acidic amino acid and N-terminal amino acid should be Basic amino acid. Then how many possible peptapeptides possible ?
\qn{5} How many of the following are Addition as well as Co-polymers ?\\
Teflon, Acrilan, Terylene, Nylon-6,6, Bakelite, Buna -N, Neoprene, PHBV, Glyptal, Novolac.
\qn{6} Number of antibiotic drugs from the following\\
Veronal, Equanil, Morphine, Pencillin, Ofloxacin, Aspirin, Valium, Chloramphenicol, Bithionol.
\qn{7} Flask contains A and B. Both decompose by first order kinetics. The half-lives for A and B are 3 sec and 1.8 sec respectively. If the concentrations of A and B are equal initially, the time required for the concentration of A to be four times that B (in sec)
\qn{8} The graph of $\log\dfrac xm$ vs $\log p$ for an adsorption process is a straight line inclined at an angle of $45^\circ$with intercept equals to 0.903. The mass of gas adsorbed per unit mass of adsorbent at the pressure of 0.5 atm is \underline{\hspace{1.5cm}}
\partsec{3}{MATHEMATICS}
\sect{Section-I(i)}
\qn{1} If $f(x)+f(y)=\dfrac1x+\dfrac1y,\ \forall x,y\in R-\{0\}$ and $\displaystyle\int_2^3\dfrac{3(f(x))^5-f(x)}{1-(f(x))^4}dx=\dfrac12\log\dfrac{2^\alpha}{3^\beta}$ then
\begin{qoptions}
\opt{A}{$\alpha>\beta$}
\opt{B}{$\beta$ is prime}
\opt{C}{$\alpha<2\beta$}
\opt{D}{$(\alpha+\beta)$ is prime}
\end{qoptions}
\qn{2} Let\\
$f(x)=\dfrac1{x+1}+\dfrac{2x}{(x+1)(x+2)}+\dfrac{3x^2}{(x+1)(x+2)(x+3)}+\ldots+\dfrac{nx^{n-1}}{(x+1)(x+2)(x+3)\ldots(x+n)}$ then
\begin{qoptions}
\opt{A}{$f(x)=\dfrac{x}{1+x}-\dfrac{x^n}{(x+1)(x+2)\ldots(x+n)}$}
\opt{B}{$f(x)=1-\dfrac{x^n}{(x+1)(x+2)\ldots(x+n)}$}
\opt{C}{$f'(x)=\left(-\dfrac{x^n}{(x+1)(x+2)\ldots(x+n)}\right)\left(\displaystyle\sum_{r=1}^{n}\dfrac{r}{x+r}\right)$}
\opt{D}{$f'(x)=\left(-\dfrac{x^{n-1}}{(x+1)(x+2)\ldots(x+n)}\right)\left(\displaystyle\sum_{r=1}^{n}\dfrac{r}{x+r}\right)$}
\end{qoptions}
\qn{3} Let $f(x)=\sin x+ax+b$. Then equation $f(x)=0$ has
\begin{qoptions}
\opt{A}{Only one real root which is positive if $a>1,\ b<0$}
\opt{B}{Only one real root which is negative if $a>1,\ b>0$}
\opt{C}{Only one real root which is negative if $a<-1,\ b<0$}
\opt{D}{Only one real root which is positive if $a<-1,\ b<0$}
\end{qoptions}
\qn{4} Let $f(x)=x+\displaystyle\int_{-1}^{1}\left(x-tx^2\right)f(t)\,dt$. Then which of the following alternative(s) is/are correct?
\begin{qoptions}
\opt{A}{$\displaystyle\lim_{x\to0}f(x)=2$}
\opt{B}{$f'\left(\dfrac{-1}3\right)=1$}
\opt{C}{$f$ is continuous and derivable on $R$.}
\opt{D}{Maximum value of $f(x)$ does not exist}
\end{qoptions}
\qn{5} Let $f:R\to R,\ y=f(x),\ f(0)=0,\ f'(x)>0$ and $f''(x)>0$. If three points $A(\alpha,f(\alpha)),B(\beta,f(\beta)),C(\gamma,f(\gamma))$ lie on $y=f(x)$ such that $0<\alpha<\beta<\gamma$, then which is/are true?
\begin{qoptions}
\opt{A}{$\alpha f(\beta)>\beta f(\alpha)$}
\opt{B}{$\dfrac{f(\alpha)+f(\beta)}{2}>f\left(\dfrac{\alpha+\beta}2\right)$}
\opt{C}{$\gamma f(\gamma+\beta-\alpha)>(\gamma+\beta-\alpha)f(\gamma)$}
\opt{D}{$\dfrac{f(\alpha)+f(\beta)}{2}<f\left(\dfrac{\alpha+\beta}2\right)$}
\end{qoptions}
\qn{6} If $\alpha,\beta\in R,\beta\ne0,n\in N$ and $\displaystyle\lim_{n\to\infty}\dfrac{\sqrt[n^2]{(1!)(2!)(3!)\ldots(n!)}}{n^\alpha}=\beta$ then
\begin{qoptions}
\opt{A}{$\alpha=\dfrac12$}
\opt{B}{$\beta=e^{\frac{-3}4}$}
\opt{C}{$\alpha=1$}
\opt{D}{$\beta=\dfrac1{\sqrt e}$}
\end{qoptions}
\qn{7} Which of the following function(s) $f(x)$ cannot exist such that?
\begin{qoptions}
\opt{A}{$f''(x)>0\ \forall x\in R,\ f'(0)=1$ and $f'(1)=1$}
\opt{B}{$f''(x)>0\ \forall x\in R,\ f'(0)=1$ and $f'(1)=0$}
\opt{C}{$f''(x)>0\ \forall x\in R,\ f'(0)=1$ and $\displaystyle\lim_{x\to\infty}f(x)=2017$}
\opt{D}{$f''(x)>0\ \forall x\in R,\ f'(0)=1$ and $f(x)\leqslant1\ \forall x<0$}
\end{qoptions}
\qn{8} Let $g(x)=\dfrac{\cos\pi x}{x(1-x)}$, then which of the following is/are correct?
\begin{qoptions}
\opt{A}{$g''\left(\dfrac12\right)=0$}
\opt{B}{$\displaystyle\int_0^1g(x)\,dx=0$}
\opt{C}{$g'\left(\dfrac12\right)=-\dfrac12$}
\opt{D}{$g(x)<0\ \forall x\in\left(\dfrac12,1\right)$}
\end{qoptions}
\sect{Section-I(ii)}
\paraq{9}{10}
Define a function $f:(0,\infty)\to B$\\
$f(x)=x\displaystyle\int_1^x\dfrac{e^t\,dt}{t}-e^x$ then
\qn{9} Over the interval $(0,\infty)$
\begin{qoptions}
\opt{A}{$f(x)$ has a point of local minima and a point of local maximum}
\opt{B}{$f(x)$ has a point of local minima but no point of local maximum}
\opt{C}{$f(x)$ has no point of local minima and has no point of local maximum}
\opt{D}{$f(x)$ has a point of local maxima but no point of local minima}
\end{qoptions}
\qn{10} Which of the following options are correct
\begin{qoptions}
\opt{A}{If $B=[-e,\infty)$ then $f(x)$ is onto function}
\opt{B}{If $B=[-e,\infty)$ then $f(x)$ is into function}
\opt{C}{In $x\in(0,1]$ equation $f(x)=e^x(x\ln x-1)$ has atleast 2 solutions}
\opt{D}{In $x\in(0,1]$ equation $f(x)=e^x(x\ln x-1)$ has exactly one solution}
\end{qoptions}
\paraq{11}{12}
Consider expression\\
$L=\sqrt{2a^2-8a+10}+\sqrt{b^2-6b+10}+\sqrt{2a^2+b^2-2ab}$ where $a,b\in R^+$
\qn{11} Minimum value of $L$ is greater than
\begin{qoptions}
\opt{A}{3}
\opt{B}{4}
\opt{C}{5}
\opt{D}{6}
\end{qoptions}
\qn{12} When $L$ is minimum then
\begin{qoptions}
\opt{A}{$3a+2b=10$}
\opt{B}{$3a+2b=11$}
\opt{C}{$a=2$}
\opt{D}{$b=5/2$}
\end{qoptions}
\sect{Section-III}
\qn{1} If the integral part of $\displaystyle\sum_{n=2}^{9999}\dfrac1{\sqrt n}$ is equal to $k$ then last digit of $k$ is \underline{\hspace{1cm}}
\qn{2} Let $f(x)=x^3-\dfrac{3x^2}2+x+\dfrac14$, then the value of $\left\{\displaystyle\int_{1/4}^{3/4}f(f(x))\,dx\right\}^{-1}$ is
\qn{3} If $k\cos^2x-k\cos x+2\geqslant0\ \forall x\in R$ then number of possible negative integral values of $k$ is :
\qn{4} If the value of $\displaystyle\int_0^2x\left\{\dfrac{d^2}{dx^2}\left(4-x^2\right)^4\right\}dx$ is $N$, then the value of $\dfrac N{32}$ is
\qn{5} If $\alpha,\beta,\gamma$ are the roots of cubic equation $2x^3-4x^2+x-8=0$ and $a=\left|5\left(\dfrac1{\alpha\beta+\gamma-1}+\dfrac1{\alpha\gamma+\beta-1}+\dfrac1{\beta\gamma+\alpha-1}\right)\right|$ then $\dfrac{10}a$ is equal to
\qn{6} A differentiable function $y=f(x)$ satisfy $\displaystyle\int_0^x(x-t+1)f(t)\,dt=x^4+x^2\ \forall x\geqslant0$ if $f(0)=26$ then minimum value of $\dfrac{f(x)}2$ is \underline{\hspace{1cm}}
\qn{7} Let $y$ be the solution of the differential equation $\left(1-x^2\right)dy=\left(xy+\left(x^3+2\right)\sqrt{1-x^2}\right)dx,\ -1<x<1$ and $y(0)=0$ if $\displaystyle\int_{-\frac12}^{\frac12}\sqrt{1-x^2}\,y(x)\,dx=k$ then $\left[\dfrac1{50k}\right]$ is equal to : ( where $[.]$ is greatest integer function)
\qn{8} Let $y=y(x)$ be the solution of the differential equation $\dfrac{dy}{dx}+\dfrac{\sqrt2y}{2\cos^4x-\cos2x}=xe^{\tan^{-1}\left(\sqrt2\cot2x\right)},\ 0<x<\pi/2$ with $y\left(\dfrac\pi4\right)=\dfrac{\pi^2}{32}$. If $y\left(\dfrac\pi3\right)=\dfrac{\pi^2}{18}e^{-\tan^{-1}(\alpha)}$ then the value of $3\alpha^2$ is equal to \underline{\hspace{1cm}}
\end{multicols}
\end{document}
